\documentclass[a4paper,12pt]{article}
\usepackage[T2A]{fontenc}
\usepackage[utf8]{inputenc}
\usepackage[russian]{babel}
\usepackage{amsmath,amssymb,mathtools}
\usepackage{helvet}
\renewcommand{\familydefault}{\sfdefault}
\usepackage{icomma}
\usepackage[a4paper,margin=20mm]{geometry}
\usepackage{microtype}
\usepackage{tikz}
\usetikzlibrary{arrows.meta,positioning,shapes.geometric,calc}
\usepackage{tabularx,booktabs,array}
\usepackage{enumitem}
\usepackage{hyperref}
\usepackage{parskip}
\usepackage{fancyhdr}
\usepackage{setspace}
\usepackage{xcolor}

\definecolor{accent}{HTML}{2C6E6F}
\definecolor{darkaccent}{HTML}{214F50}
\definecolor{textgray}{HTML}{2F3333}
\definecolor{softgray}{HTML}{6B7373}
\definecolor{linegray}{HTML}{D8DEDE}

\hypersetup{colorlinks=true,linkcolor=darkaccent,urlcolor=darkaccent}
\setlength{\headheight}{25pt}
\onehalfspacing
\setlength{\parindent}{0pt}
\setlist[itemize]{leftmargin=1.4em,itemsep=0.25em,topsep=0.3em}
\setlist[enumerate]{leftmargin=1.7em,itemsep=0.55em,topsep=0.4em}
\pagestyle{fancy}
\fancyhf{}
\fancyhead[L]{\small\color{softgray} Обыкновенные дроби}
\fancyhead[R]{\small\color{softgray} 06.10.26}
\fancyfoot[C]{\small\color{softgray}\thepage}
\renewcommand{\headrulewidth}{0.3pt}
\renewcommand{\headrule}{\hbox to\headwidth{\color{linegray}\leaders\hrule height \headrulewidth\hfill}}

\newcommand{\sectiontitle}[1]{%
  \vspace{0.7em}
  {\large\bfseries\color{darkaccent}#1}\par
  \vspace{0.2em}\color{linegray}\hrule\color{textgray}\vspace{0.55em}
}
\newcommand{\keyidea}[1]{\textbf{Ключевая идея.} #1}
\newcommand{\warning}[1]{\textbf{Важно.} #1}
\newcommand{\ex}[1]{\textbf{Пример.} #1}

\begin{document}
\color{textgray}

\begin{titlepage}
\thispagestyle{empty}
\centering
\vspace*{26mm}
{\small\color{softgray} МАТЕМАТИКА}\par
\vspace{8mm}
{\Huge\bfseries\color{darkaccent} Чек-лист}\par
\vspace{5mm}
{\LARGE\bfseries Обыкновенные дроби}\par
\vspace{2mm}
{\large общий знаменатель, вычисления, уравнения и задачи}\par
\vspace{15mm}
{\large 06.10.26}\par
\vfill
{\large\bfseries Лёвин Артём Александрович}\par
\vspace{2mm}
{\normalsize эксклюзивно для Нади Клименко}\par
\vspace{17mm}
\begin{tikzpicture}[remember picture,overlay]
  \draw[accent,line width=1.15pt]
  ($(current page.south west)+(18mm,20mm)$)
  .. controls ($(current page.south west)+(62mm,31mm)$)
  and ($(current page.south east)+(-68mm,10mm)$)
  .. ($(current page.south east)+(-18mm,20mm)$);
\end{tikzpicture}
\end{titlepage}

\sectiontitle{1. Общий знаменатель: что именно нужно уметь}

\keyidea{Чтобы сравнивать, складывать или вычитать дроби с разными знаменателями, приведите их к общему знаменателю. Удобнее всего использовать наименьший общий знаменатель, то есть НОК знаменателей.}

Для дробей \(\frac{a}{m}\) и \(\frac{b}{n}\):
\[
L=\mathop{\text{НОК}}(m,n),\qquad
\frac{a}{m}=\frac{a\cdot(L/m)}{L},\qquad
\frac{b}{n}=\frac{b\cdot(L/n)}{L}.
\]

\ex{Сравним \(\frac{26}{55}\) и \(\frac{16}{35}\).}
\[
55=5\cdot 11,\qquad 35=5\cdot 7,\qquad \mathop{\text{НОК}}(55,35)=385.
\]
\[
\frac{26}{55}=\frac{182}{385},\qquad
\frac{16}{35}=\frac{176}{385}.
\]
Следовательно,
\[
\boxed{\frac{26}{55}>\frac{16}{35}}.
\]

\warning{Дополнительный множитель лучше находить по разложению НОК, а не длинным делением. Для знаменателя \(55=5\cdot11\) в НОК \(5\cdot7\cdot11\) не хватает множителя \(7\).}

\sectiontitle{2. Сначала упростить, потом считать}

Перед поиском общего знаменателя осмотрите выражение. Часто его можно сделать заметно короче:
\begin{itemize}
  \item сократить отдельную дробь;
  \item сложить дроби с уже одинаковыми знаменателями;
  \item убрать скобки после знака «+»;
  \item выполнить удобное действие внутри скобок;
  \item только затем искать общий знаменатель для оставшихся дробей.
\end{itemize}

\ex{Вычислим}
\[
\frac{23}{24}-\left(\frac16+\frac14\right).
\]
Внутри скобок сразу удобно взять знаменатель \(12\):
\[
\frac16+\frac14=\frac2{12}+\frac3{12}=\frac5{12}=\frac{10}{24}.
\]
Тогда
\[
\frac{23}{24}-\frac{10}{24}=\boxed{\frac{13}{24}}.
\]

\warning{Сокращать «по диагонали» через знак сложения или вычитания нельзя. Сокращение работает только для общих множителей числителя и знаменателя одной дроби или произведения.}

\sectiontitle{3. Десятичные дроби переводим в обыкновенные}

Десятичную дробь сначала запишите со знаменателем \(10\), \(100\), \(1000\) и т.~д., затем обязательно сократите.

\begin{center}
\renewcommand{\arraystretch}{1.35}
\begin{tabular}{ccccc}
\toprule
\(0{,}2\) & \(0{,}25\) & \(0{,}4\) & \(0{,}75\) & \(0{,}125\) \\
\midrule
\(\frac15\) & \(\frac14\) & \(\frac25\) & \(\frac34\) & \(\frac18\) \\
\bottomrule
\end{tabular}
\end{center}

\ex{
\[
0{,}2+\frac17=\frac15+\frac17=\frac7{35}+\frac5{35}=\boxed{\frac{12}{35}}.
\]
}

\sectiontitle{4. Сравнение без громоздких вычислений}

Иногда полезно сначала сравнить дроби с опорными значениями \(\frac12\), \(\frac34\) или \(1\). Это быстро показывает порядок величин или хотя бы резко сужает сравнение.

Например,
\[
\frac7{12}<\frac34,\qquad
\frac{11}{13}>\frac34,
\]
поэтому сразу ясно, что \(\frac7{12}<\frac{11}{13}\).

Чтобы сравнить \(\frac{11}{13}\) и \(\frac{57}{60}\), удобно посмотреть, сколько не хватает до единицы:
\[
1-\frac{11}{13}=\frac2{13},\qquad
1-\frac{57}{60}=\frac1{20}.
\]
Так как \(\frac1{20}<\frac2{13}\), дробь \(\frac{57}{60}\) ближе к единице и потому больше. Итак,
\[
\boxed{\frac7{12}<\frac{11}{13}<\frac{57}{60}}.
\]

\warning{Опорные дроби дают быстрый ориентир. Если две дроби оказались по одну сторону от ориентира и порядок не очевиден, используйте точное сравнение: общий знаменатель или перекрёстное умножение.}

\sectiontitle{5. Уравнения с дробями: сначала изолируйте нужный объект}

\keyidea{Не начинайте с общего знаменателя автоматически. Сначала упростите структуру уравнения и выразите неизвестное или целый блок, в котором оно находится.}

\ex{Решим}
\[
t-\frac{11}{18}=\frac{11}{12}-\frac59.
\]
Сначала переносим \(-\frac{11}{18}\) вправо:
\[
t=\frac{11}{12}-\frac59+\frac{11}{18}.
\]
Теперь общий знаменатель \(36\):
\[
t=\frac{33-20+22}{36}=\boxed{\frac{35}{36}}.
\]

\ex{Решим}
\[
\frac45-\left(\frac9{10}-z\right)=\frac15.
\]
Сначала выражаем целую скобку:
\[
\frac9{10}-z=\frac45-\frac15=\frac35.
\]
Затем
\[
z=\frac9{10}-\frac35=\boxed{\frac3{10}}.
\]

\sectiontitle{6. Текстовые задачи: дроби остаются теми же действиями}

\textbf{Периметр.} Если одна сторона неизвестна, обозначьте её \(x\) и используйте сумму сторон.

\ex{Пусть \(P=\frac{17}{20}\) м, \(AB=\frac{17}{50}\) м, а \(BC\) на \(\frac3{50}\) м длиннее \(AB\). Тогда}
\[
BC=\frac{20}{50}=\frac25,
\]
\[
\frac{17}{50}+\frac25+x=\frac{17}{20},
\qquad
x=\boxed{\frac{11}{100}\text{ м}}.
\]

\textbf{Совместная работа.} Если весь заказ принимаем за \(1\), а мастер выполняет его за \(T\) часов, его производительность равна \(\frac1T\) заказа в час. За \(t\) часов он выполнит \(\frac{t}{T}\) заказа.

\ex{Первый мастер выполняет заказ за \(14\) ч, второй --- за \(18\) ч. За \(5\) и \(7\) часов соответственно они выполнят}
\[
\frac5{14}+\frac7{18}=\frac{45}{126}+\frac{49}{126}=\boxed{\frac{47}{63}}
\]
заказа.

\textbf{Догонка.} Если объекты движутся в одном направлении, расстояние между ними сокращается со скоростью \(v_{\text{быстр}}-v_{\text{медл}}\).

\ex{Начальный разрыв \(10\) км, скорости \(1{,}7\) и \(1{,}2\) км/мин:}
\[
(1{,}7-1{,}2)t=10,\qquad 0{,}5t=10,\qquad \boxed{t=20\text{ мин}}.
\]

\sectiontitle{7. Мини-памятка перед ответом}

\begin{itemize}
  \item Сокращена ли каждая дробь, которую можно сократить?
  \item Можно ли упростить выражение до поиска общего знаменателя?
  \item НОК действительно наименьший и удобный?
  \item Дополнительный множитель умножает и числитель, и знаменатель?
  \item В уравнении сначала изолирован неизвестный или нужный блок?
  \item В текстовой задаче ясно, что обозначено через \(x\), и сохранены единицы измерения?
  \item Финальная дробь сокращена?
\end{itemize}

\clearpage
\thispagestyle{fancy}
\begin{center}
{\Large\bfseries\color{darkaccent} Алгоритм: рациональная работа с дробным выражением}
\end{center}
\vspace{8mm}

\begin{center}
\begin{tikzpicture}[
  node distance=11mm,
  startstop/.style={ellipse,draw=accent,line width=0.9pt,minimum width=42mm,minimum height=10mm,align=center,font=\small},
  process/.style={rounded corners=2mm,draw=accent,line width=0.9pt,minimum width=92mm,minimum height=12mm,text width=80mm,align=center,font=\small},
  decision/.style={diamond,draw=accent,line width=0.9pt,aspect=2.4,text width=38mm,align=center,font=\small,inner sep=1.2mm},
  arrow/.style={-{Stealth[length=2.6mm]},line width=0.9pt,color=darkaccent}
]
\node[startstop] (start) {Получено выражение с дробями};
\node[decision,below=of start] (simple) {Есть что упростить до НОК?};
\node[process,below left=14mm and -8mm of simple] (reduce) {Сократите отдельные дроби, выполните очевидные действия, уберите лишние скобки};
\node[process,below=34mm of simple] (lcm) {Найдите НОК оставшихся знаменателей и дополнительные множители};
\node[process,below=of lcm] (convert) {Приведите дроби к общему знаменателю};
\node[process,below=of convert] (calc) {Выполните сложение или вычитание числителей};
\node[decision,below=of calc] (canred) {Итоговая дробь сокращается?};
\node[process,below left=14mm and -8mm of canred] (finalred) {Сократите дробь};
\node[startstop,below=34mm of canred] (finish) {Запишите окончательный ответ};

\draw[arrow] (start) -- (simple);
\draw[arrow] (simple) -- node[above left,font=\small]{да} (reduce);
\draw[arrow] (reduce.south) |- (lcm.west);
\draw[arrow] (simple) -- node[right,font=\small]{нет} (lcm);
\draw[arrow] (lcm) -- (convert);
\draw[arrow] (convert) -- (calc);
\draw[arrow] (calc) -- (canred);
\draw[arrow] (canred) -- node[above left,font=\small]{да} (finalred);
\draw[arrow] (finalred.south) |- (finish.west);
\draw[arrow] (canred) -- node[right,font=\small]{нет} (finish);
\end{tikzpicture}
\end{center}

\clearpage
\sectiontitle{Тренировочный блок --- Путь А}

\textbf{10 заданий.} Решайте по порядку: от базовой техники к задачам, где нужно выбрать удобную модель.

\begin{enumerate}
  \item Сравните дроби \(\dfrac7{18}\) и \(\dfrac5{12}\).

  \item Расположите по возрастанию дроби
  \[
  \frac{11}{16},\qquad \frac{17}{24},\qquad \frac{23}{40}.
  \]
  Постарайтесь сначала использовать сравнение с удобными опорными дробями.

  \item Вычислите:
  \[
  \frac5{24}+\frac7{36}.
  \]

  \item Вычислите рациональным способом:
  \[
  \frac{23}{24}-\left(\frac16+\frac14\right).
  \]

  \item Представьте десятичную дробь в виде несократимой обыкновенной и вычислите:
  \[
  0{,}25+\frac7{12}.
  \]

  \item Решите уравнение:
  \[
  t-\frac{11}{18}=\frac{11}{12}-\frac59.
  \]

  \item Решите уравнение, сначала выразив скобку:
  \[
  \frac45-\left(\frac9{10}-z\right)=\frac15.
  \]

  \item Периметр треугольника равен \(\dfrac{23}{20}\) м. Одна сторона равна \(\dfrac7{20}\) м, вторая на \(\dfrac1{10}\) м длиннее первой. Найдите третью сторону.

  \item Первый мастер выполняет весь заказ за \(15\) часов, второй --- за \(20\) часов. Какую часть заказа они выполнят вместе, если первый работает \(6\) часов, а второй --- \(8\) часов?

  \item Два автомобиля движутся в одном направлении. В начале движения более быстрый автомобиль находится на \(7{,}5\) км позади другого. Их скорости равны \(1{,}8\) км/мин и \(1{,}3\) км/мин. Через сколько минут быстрый автомобиль догонит медленный?
\end{enumerate}

\clearpage
\sectiontitle{Ответы}

\renewcommand{\arraystretch}{1.45}
\begin{tabularx}{\linewidth}{>{\bfseries}c X}
\toprule
№ & Ответ \\
\midrule
1 & \(\dfrac7{18}<\dfrac5{12}\) \\
2 & \(\dfrac{23}{40}<\dfrac{11}{16}<\dfrac{17}{24}\) \\
3 & \(\dfrac{29}{72}\) \\
4 & \(\dfrac{13}{24}\) \\
5 & \(\dfrac56\) \\
6 & \(t=\dfrac{35}{36}\) \\
7 & \(z=\dfrac3{10}\) \\
8 & \(\dfrac7{20}\text{ м}\) \\
9 & \(\dfrac45\) заказа \\
10 & \(15\text{ мин}\) \\
\bottomrule
\end{tabularx}

\vfill
{\small\color{softgray} Лёвин Артём Александрович эксклюзивно для Нади Клименко}

\end{document}
